% section_1.tex — Background and Setup

\section{Background and Setup}
\label{sec:background}

Novel View Synthesis (NVS) is an active area of research in 3D computer
graphics. The goal is to reconstruct a 3D representation of a scene from
a sparse collection of photographs taken from multiple camera angles.
Optimization-based approaches that fit either differentiable neural
fields~\citep{mildenhall2020nerf,mueller2022instantngp} or Gaussian
primitives~\citep{kerbl20233dgs} have advanced the field substantially
in both rendering quality and inference speed; we refer readers
to~\citet{tewari2022nrsurvey} for a comprehensive survey.

Defining ``good quality'' for a 3D rendering method is open-ended: the
final product is a visual artifact meant to be perceived by human eyes.
The literature predominantly relies on image-based metrics --- PSNR,
MS-SSIM, and LPIPS~\citep{zhang2018lpips} --- which compare reconstructed
views against ground-truth photographs held out from training. As
numerical distances these metrics are informative, but they have two
limitations in the 3D-rendering context:
\begin{enumerate}
    \item Two decades of image-quality work shows that good PSNR or
    MS-SSIM scores do not necessarily translate to a good perceptual
    experience~\citep{wang2009mse}. LPIPS partially addresses this with
    learned deep features but remains a proxy for direct viewing and is
    expensive to compute.
    \item Image-based metrics evaluate \emph{static} images at a small,
    fixed set of held-out camera poses, while users actually
    \emph{interact} with 3D scenes --- moving the camera, zooming, and
    viewing continuously from arbitrary angles. Free navigation
    introduces failure modes (popping, floater drift, smoothness
    artifacts at oblique angles) that no static-image metric can detect.
\end{enumerate}

To address this gap, we develop an interactive streaming harness for
human evaluation of real-time 3D renderers, built around the
philosophy that rendering quality should ultimately be measured by how
participants perceive it under interactive viewing. A researcher
supplies one or two renderers and a study configuration; the system
streams to participants' browsers, collects per-aspect ratings (and,
in two-method studies, an overall A/B preference with aspect
citations), and returns the raw per-participant data.

The system's inputs are a researcher-provided rendering algorithm and
its fitted 3D model, together with a study configuration (scenes,
conditions, protocol). Its outputs are the data collected from each
participant's session: (a)~\emph{subject ratings} produced through the
rating UI --- per-aspect 0--100 sliders, A/B preferences, and aspect
citations; (b)~\emph{free-text feedback} entered in the post-session
notes field; and (c)~\emph{mouse-movement trajectories} logged over
the participant's viewing window. The first two are the immediate
data products of the rating task; the third lets a researcher
reconstruct \emph{which views} of the scene each participant chose to
inspect before issuing a rating. Figure~\ref{fig:io} gives a
high-level view of how the system sits between the researcher and the
participant.

\begin{figure}[t]
  \centering
  \includegraphics[width=\linewidth]{io_diagram}
  \caption{Conceptual view of system inputs and outputs. The researcher
  provides a rendering algorithm and a fitted 3D model; the system
  streams the resulting renderings to a participant's browser-based
  viewer and collects three categories of response data --- subject
  ratings, free-text feedback, and mouse-movement trajectories ---
  for the researcher to analyze.}
  \label{fig:io}
\end{figure}

The three \textit{main goals} are:
\begin{enumerate}
    \item Measure a dimension of rendering quality that image-based
    metrics on held-out static frames cannot capture: the participant's
    \emph{interactive perceptual experience} under continuous,
    free-roaming viewing. The collected data is intended as citable
    evidence a researcher can use to support claims about how their
    method actually behaves in interactive use, complementing rather
    than competing with the offline numbers already reported in the
    literature.
    \item Offer a plug-and-play framework where a researcher can connect
    their renderer and trained model and run a human study with minimal
    infrastructure work.
    \item Support both single-method subjective studies and two-method
    comparison studies, covering the two most common shapes of perceptual
    question a 3D-graphics researcher might ask.
\end{enumerate}
The three explicit \textit{non-goals} are:
\begin{enumerate}
    \item Replace image-based numerical metrics --- those metrics are
    cheap to compute, can be used as training losses, and remain the
    right tool for in-the-loop comparisons during method development.
    \item Modify the underlying rendering algorithms --- we provide a
    platform for evaluation, not for method improvement.
    \item Provide statistical analysis or visualization of the
    collected data --- the harness delivers raw per-participant
    response data; how that data is summarized, correlated against
    offline metrics, or visualized is the researcher's call.
\end{enumerate}
Two practical \textit{constraints} shape the system's design:
\begin{enumerate}
    \item \textit{Participants evaluate from their own machines},
    which we cannot assume to have GPUs capable of running an NVS
    renderer. The renderer must therefore run on a remote
    GPU-equipped instance, with the rendered output delivered to the
    participant's browser over the network. We host the active
    renderer(s) and supporting services on an AWS L4 GPU instance.
    \item \textit{The framework should be agnostic to the underlying
    rendering algorithm.} Different NVS methods use different
    rendering stacks --- CUDA-accelerated hash grids for i-NGP, a
    custom rasterizer for Nexels, the \texttt{gsplat} differentiable
    rasterizer for 3DGS --- and we want to support not just these
    three but whatever NVS method a future researcher brings. A
    framework that asked each renderer to ship its native output
    format to the browser would have to keep pace with every method's
    evolving toolchain. Streaming a common video format (H.264)
    abstracts that diversity away: every renderer emits the same
    output regardless of its internal representation.
\end{enumerate}
These constraints together motivate the streaming architecture
described in Section~\ref{sec:approach}.

To exercise the harness, we run two case studies of the kind a graphics
researcher might reach for when validating their own model. The first is
a single-viewpoint study --- participants view one rendering at a time
and rate it. The second is a two-viewpoint study --- participants compare
two NVS models side by side and select an overall preference along with
the aspects that drove that choice. The focus is on the system rather
than the findings of either case study; provided the resulting data are
coherent and non-degenerate, the harness has done its job.

A good evaluation harness must succeed on three independent axes.
\emph{System-side}: the streaming pipeline has to work end-to-end for a
diverse set of NVS methods. \emph{Researcher-side}: the data has to be
informative and analyzable on the researcher's own terms.
\emph{Participant-side}: each session has to feel smooth, clear, and
easy. We make all three criteria precise in
Section~\ref{sec:results}. The crux is twofold, and the methodological
half is the larger contribution. \emph{Methodologically}, we have to
design a questionnaire and a study protocol that meaningfully capture
interactive perceptual experience --- a design problem with no
off-the-shelf answer in the 3D-rendering literature, where every
existing instrument either reduces to an image-metric proxy or asks the
participant to rate a static frame. \emph{Infrastructurally}, we have
to deliver real-time renderings of multiple high-end NVS methods to a
remote browser at participant-imperceptible latency, on a single L4
GPU, with a renderer interface generic enough that adding a new method
does not require modifying the study logic.
Section~\ref{sec:approach} describes both;
Section~\ref{sec:results} reports the two case studies.
